% Options for packages loaded elsewhere
% Options for packages loaded elsewhere
\PassOptionsToPackage{unicode}{hyperref}
\PassOptionsToPackage{hyphens}{url}
\PassOptionsToPackage{dvipsnames,svgnames,x11names}{xcolor}
%
\documentclass[
  letterpaper,
  DIV=11,
  numbers=noendperiod]{scrartcl}
\usepackage{xcolor}
\usepackage{amsmath,amssymb}
\setcounter{secnumdepth}{-\maxdimen} % remove section numbering
\usepackage{iftex}
\ifPDFTeX
  \usepackage[T1]{fontenc}
  \usepackage[utf8]{inputenc}
  \usepackage{textcomp} % provide euro and other symbols
\else % if luatex or xetex
  \usepackage{unicode-math} % this also loads fontspec
  \defaultfontfeatures{Scale=MatchLowercase}
  \defaultfontfeatures[\rmfamily]{Ligatures=TeX,Scale=1}
\fi
\usepackage{lmodern}
\ifPDFTeX\else
  % xetex/luatex font selection
\fi
% Use upquote if available, for straight quotes in verbatim environments
\IfFileExists{upquote.sty}{\usepackage{upquote}}{}
\IfFileExists{microtype.sty}{% use microtype if available
  \usepackage[]{microtype}
  \UseMicrotypeSet[protrusion]{basicmath} % disable protrusion for tt fonts
}{}
\makeatletter
\@ifundefined{KOMAClassName}{% if non-KOMA class
  \IfFileExists{parskip.sty}{%
    \usepackage{parskip}
  }{% else
    \setlength{\parindent}{0pt}
    \setlength{\parskip}{6pt plus 2pt minus 1pt}}
}{% if KOMA class
  \KOMAoptions{parskip=half}}
\makeatother
% Make \paragraph and \subparagraph free-standing
\makeatletter
\ifx\paragraph\undefined\else
  \let\oldparagraph\paragraph
  \renewcommand{\paragraph}{
    \@ifstar
      \xxxParagraphStar
      \xxxParagraphNoStar
  }
  \newcommand{\xxxParagraphStar}[1]{\oldparagraph*{#1}\mbox{}}
  \newcommand{\xxxParagraphNoStar}[1]{\oldparagraph{#1}\mbox{}}
\fi
\ifx\subparagraph\undefined\else
  \let\oldsubparagraph\subparagraph
  \renewcommand{\subparagraph}{
    \@ifstar
      \xxxSubParagraphStar
      \xxxSubParagraphNoStar
  }
  \newcommand{\xxxSubParagraphStar}[1]{\oldsubparagraph*{#1}\mbox{}}
  \newcommand{\xxxSubParagraphNoStar}[1]{\oldsubparagraph{#1}\mbox{}}
\fi
\makeatother


\usepackage{longtable,booktabs,array}
\usepackage{calc} % for calculating minipage widths
% Correct order of tables after \paragraph or \subparagraph
\usepackage{etoolbox}
\makeatletter
\patchcmd\longtable{\par}{\if@noskipsec\mbox{}\fi\par}{}{}
\makeatother
% Allow footnotes in longtable head/foot
\IfFileExists{footnotehyper.sty}{\usepackage{footnotehyper}}{\usepackage{footnote}}
\makesavenoteenv{longtable}
\usepackage{graphicx}
\makeatletter
\newsavebox\pandoc@box
\newcommand*\pandocbounded[1]{% scales image to fit in text height/width
  \sbox\pandoc@box{#1}%
  \Gscale@div\@tempa{\textheight}{\dimexpr\ht\pandoc@box+\dp\pandoc@box\relax}%
  \Gscale@div\@tempb{\linewidth}{\wd\pandoc@box}%
  \ifdim\@tempb\p@<\@tempa\p@\let\@tempa\@tempb\fi% select the smaller of both
  \ifdim\@tempa\p@<\p@\scalebox{\@tempa}{\usebox\pandoc@box}%
  \else\usebox{\pandoc@box}%
  \fi%
}
% Set default figure placement to htbp
\def\fps@figure{htbp}
\makeatother





\setlength{\emergencystretch}{3em} % prevent overfull lines

\providecommand{\tightlist}{%
  \setlength{\itemsep}{0pt}\setlength{\parskip}{0pt}}



 


\KOMAoption{captions}{tableheading}
\makeatletter
\@ifpackageloaded{caption}{}{\usepackage{caption}}
\AtBeginDocument{%
\ifdefined\contentsname
  \renewcommand*\contentsname{Table of contents}
\else
  \newcommand\contentsname{Table of contents}
\fi
\ifdefined\listfigurename
  \renewcommand*\listfigurename{List of Figures}
\else
  \newcommand\listfigurename{List of Figures}
\fi
\ifdefined\listtablename
  \renewcommand*\listtablename{List of Tables}
\else
  \newcommand\listtablename{List of Tables}
\fi
\ifdefined\figurename
  \renewcommand*\figurename{Figure}
\else
  \newcommand\figurename{Figure}
\fi
\ifdefined\tablename
  \renewcommand*\tablename{Table}
\else
  \newcommand\tablename{Table}
\fi
}
\@ifpackageloaded{float}{}{\usepackage{float}}
\floatstyle{ruled}
\@ifundefined{c@chapter}{\newfloat{codelisting}{h}{lop}}{\newfloat{codelisting}{h}{lop}[chapter]}
\floatname{codelisting}{Listing}
\newcommand*\listoflistings{\listof{codelisting}{List of Listings}}
\makeatother
\makeatletter
\makeatother
\makeatletter
\@ifpackageloaded{caption}{}{\usepackage{caption}}
\@ifpackageloaded{subcaption}{}{\usepackage{subcaption}}
\makeatother
\usepackage{bookmark}
\IfFileExists{xurl.sty}{\usepackage{xurl}}{} % add URL line breaks if available
\urlstyle{same}
\hypersetup{
  pdftitle={Problem Set 1},
  colorlinks=true,
  linkcolor={blue},
  filecolor={Maroon},
  citecolor={Blue},
  urlcolor={Blue},
  pdfcreator={LaTeX via pandoc}}


\title{Problem Set 1}
\author{}
\date{}
\begin{document}
\maketitle


\begin{enumerate}
\def\labelenumi{\arabic{enumi}.}
\item
  \textbf{Polio\footnote{This is a revised version of question 4.E.5
    from Cobb,
    \emph{Introduction to Design and Analysis of Experiments}, 1998.}}.
  Polio is now a very rare disease, thanks to medical research, which
  developed a vaccine, and thanks to statistical work, which proved the
  vaccine to be effective. The testing of the vaccine involved two
  different studies. One was a true experiment, with the treatment and
  placebo assigned by a coin flip, but the other study was
  observational.

  The observational study used the children who didn't get parental
  permission for the vaccine as a control group and gave the vaccine to
  everyone who got permission. In what ways do you think the treatment
  and control groups are likely to be different? How might that make it
  more challenging to identify the casual factor?
\end{enumerate}

\newpage{}

\begin{enumerate}
\def\labelenumi{\arabic{enumi}.}
\setcounter{enumi}{1}
\item
  \textbf{The Free-Choice Paradigm}. The Free-Choice Paradigm\footnote{Salti
    et al, ``Cognitive dissonance resolution is related to episodic
    memory,'\,' PLoS One, 2014.} is a term for a type of experiment
  designed to test what is termed the ``cognitive dissonance'' of a
  subject. Cognitive dissonance says that individuals strive to decrease
  the discomfort generated by conflicting cognitions (i.e.~opinions or
  thoughts) by modifying them.

  A traditional free-choice paradigm experiment has a subject do the
  following:

  \emph{Step A}. Rate items individually according to their desirability
  (e.g.~oranges, clementines, apples, bananas, etc.), say on a score
  between 1-10.

  \emph{Step B}. Choose between pairs of items that they previously
  rated similarly (``do you prefer oranges or clementines?'', ``do you
  prefer apples or bananas?'', etc.)

  \emph{Step C}. Rate items individually from the first step (again).

  Results from such experiments over the past decades show that items
  ranked highly in the 1st step get increased ratings when they are
  repeated in step C, and those ranking lowly in the 1st step get even
  lower ratings in the repetition of step C. To be concrete, consider
  the following response: for each subject, take those items rated with
  a value above 5 in step A, calculate a difference in the rankings
  between step C and step A and take an average over all such items for
  the subject. Then this response is found to be significantly greater
  than zero based on appropriate statistical tests.

  The cognitive dissonance theory explains this phenomena by the theory
  that forcing individuals to specify preferences between particular
  items in step B (``do you prefer oranges or clementines?'') forces
  them to recognize inconsistencies in their preferences and to alter
  them, hence the changes between step A and step C (i.e.~the process of
  going through step B caused the change in preferences seen in step C).

  \begin{enumerate}
  \def\labelenumii{\alph{enumii}.}
  \tightlist
  \item
    What is the nature of the treatments?
  \item
    What are the units?
  \item
    How were units assigned to treatments?
  \item
    What is the nature of the response?
  \item
    Do any of the answers to the previous questions parts make it
    difficult to assess the psychological theory that they're trying to
    test? If so, explain why.
  \end{enumerate}

  To help you answer this last question, consider the follow-up study
  proposed by a critic of this paper. The critic reruns the above
  experiment, only changing the order of the steps: step A, step C, step
  B. They find a similar value of the response (after appropriate
  statistical tests) as when running the original experiment of step A,
  step B, step C.
\end{enumerate}

\newpage{}

\begin{enumerate}
\def\labelenumi{\arabic{enumi}.}
\setcounter{enumi}{2}
\item
  \textbf{Estimation under Imbalance}. Consider the setting where we
  seek to estimate the Average Treatment Effect (ATE) using as our
  estimator the observed difference in the mean responses in the two
  groups, \(\overline{Y}_1 - \overline{Y}_0\). Let \(D_i\) track whether
  unit \(i\) receives the treatment (1) or the control (0) and let the
  set of indices in the treatment and control groups be be \(g1\) and
  \(g0\) respectively. Also let \(n_1\) and \(n_0\) be the number of
  units in each group.

  So far, we've analyzed the case where \(n_1 = n_2\), a case where the
  groups are \emph{balanced}. What happens to the estimator when the
  groups are \emph{unbalanced}?

  Determine whether or not the estimator is biased when \(n_1\) is twice
  the size of \(n_0\). Units are assigned to groups by randomly
  shuffling the list of their unique identifiers and assigning the first
  \(n_1\) to the treatment and the remaining to the control. If the
  estimator is biased, suggest an alternative bias-corrected estimator.
\end{enumerate}

\newpage{}

The last few questions will be released by 1/31/26!




\end{document}
